\section{Integrated Education Plan}
\label{sec:education}

The process of debugging---finding, reproducing and fixing bugs---is at the heart of software development as evidenced by the effort that goes into it (50-70\%~\cite{rti_2002_EconomicImpactsInadequate,goodliffe_2015_BecomingBetterProgrammer,economicimpacttesting}). I teach Operating Systems, and my office hours are almost exclusively dedicated to helping students debug. More fundamentally, at its core, debugging is a key problem-solving skill. Developers need to reason about how state is maintained, updated, and corrupted to yield to an error~\cite{debugstate}. They then need to devise the most appropriate fix. The research I propose aims to greatly aid debugging. Alas, formal computer science is education (certainly at Michigan) is almost entirely devoid of debugging material. \textit{The goal of my education plan is to integrate debugging education into outreach aimed at middle schoolers and in the undergraduate and graduate curricula.}

\noindent \textit{\textbf{Middle School Outreach.}} I will design a course aimed to teach basic principles of debugging to middle school students. I will structure the course based on recent work~\cite{debuggingtrajectories} that outlines key learning objectives for a debugging curriculum aimed at K-8 computer science education. Some examples of key learning objectives are: (1) using outcomes (outputs) to determine the presence of errors, (2) using iterative refinement to fix errors, (3) understanding the role of errors in problem solving, (4) understanding how step-by-step execution can help diagnose errors, (5) understanding the critical role of \textbf{\textit{reproducing}} bugs in debugging. 

I will design the course using the educational Microbit platform~\cite{Microbit27:online} an open-source ARM-based platform used extensively in middle school education in the UK. Microsoft develops and maintains a website, Makecode~\cite{Microsof90:online}, that can be used to program a Microbit through both graphical and textual interfaces. Therefore, the barrier to entry for using a Microbit is very reasonable (i.e., a browser and a 15\$ Microbit). 

The course will be featured on the Makecode website~\cite{Microsof90:online} and will be piloted at the Qualcomm Thinkabit Lab within the Michigan Engineering Zone~\cite{mez} (MEZ), in Detroit. MEZ is headquartered at the University of Michigan's Detroit Center, and offers resources (courses, lab space) for engineering education geared towards \textit{\textbf{underrepresented and underserved minorities}} in Southeast Michigan.

I will design the course in collaboration with a CS education expert Prof. Mark Guzdial, Thinkabit Lab~\cite{Qualcomm7:online} program coordinator Ms. Haley Hart, and the leader of the Microbit Makecode~\cite{Microsof90:online} effort at Microsoft, principal researcher, Dr. Thomas Ball (see attached letters). I will consult Michigan Center for Engineering Diversity and Outreach (see letter from the director, Mr. Lyonel Milton) when targeting the course for the Thinkabit students, which are almost exclusively underrepresented minorities. I participated in a summer camp at the Thinkabit Lab in June 2019, which gave me a chance to observe the learning environment, interact with the children, and study how instructors conduct lectures. Ms. Hart will help me pilot the course at the Thinkabit Lab, which is crucial for the success of this initiative.

\noindent \textit{\textbf{Undergraduate Education.}} I will add 2 lectures and 2 lab assignments dedicated to debugging to the contents of the undergraduate operating systems course that I teach (EECS 482). The first lecture will be about using advanced debugging tools such as record/replay, instrumentation frameworks, etc. The second will be based on real case studies from the class project itself. In this class, students build major components of an operating system, and every year, there are a number of interesting and complex bugs that they need to deal with (race conditions, deadlocks, memory errors etc.). I will walk through example bugs and how to think about solving them and focus on distilling principles for debugging OS code. Each lecture will have an accompanying lab assignment to solidify the concepts in a practical setting.

\noindent \textit{\textbf{Graduate Education.}} I teach the graduate operating systems course at Michigan (EECS 582), which I overhauled to cover topics in reliability and bug finding. I will revise this course with the most recent developments in this space. This course features a research project component where students build a prototype and write a paper about it. One of my PhD students and an undergraduate student taking this course went on to do further research with me and recently published their work at PLDI'19~\cite{huron}. In the same vein, I will use this course as a forum to recruit students to work on some of the ideas in this proposal.

\noindent \textit{\textbf{Evaluation of the Education Plan.}} Thinkabit Lab has a well-established evaluation and feedback process that I will leverage. For the undergraduate and graduate courses, I will measure success based on teaching evaluations and the number of students who join my group and/or pursue graduate school careers.


% introduce programming skills to a highly diverse student population across the entire UC San Diego campus

% Graduate Education I currently teach the graduate seminar titled “Big Data Systems and Applications”(EECS 598) [17] that involves hands-on assignments on modern memory-intensive applications such as Apache Spark and TensorFlow. Students review papers through a conference management system, leading to critical thinking and fostering better collaboration. Students also have to complete a semester-long project. I plan to integrate many of the proposed tasks as potential projects for this course in coming years. Indeed, I have had success in this manner. Prior to creating EECS 598, I revamped and taught the graduate-level Advanced Operating Systems course (EECS 582) in Winter 2016 and Fall 2016 [15, 16]. One of the course projects led to Infiniswap [101] – the precursor to this proposal.

% needs also an evaluation and success critera 

% \noindent \textbf{Educational Impact.} The proposed project has three aspects of educational impact.  
% Second, to the curriculum of \textit{EECS 582: Advanced Operating Systems} in Michigan and \textit{ECE695: Operating Systems Design and Implementation} in Purdue, graduate courses taught by the PIs, we will incorporate research paper discussion regarding IoT systems reliability.
% The resultant software from this project will be integrated into course projects. 